% Standalone problem document, generated from the adjacent README.md.
% Compile with XeLaTeX. Mathematical content is maintained in Markdown.
\documentclass[11pt,a4paper]{article}
\usepackage[fontsize=10.5pt]{scrextend}
\usepackage[margin=22mm,headheight=15pt,headsep=9mm,footskip=11mm]{geometry}
\usepackage{fontspec}
\setmainfont{texgyrepagella}[Extension=.otf,UprightFont=*-regular,BoldFont=*-bold,ItalicFont=*-italic,BoldItalicFont=*-bolditalic]
\setsansfont{texgyreheros}[Extension=.otf,UprightFont=*-regular,BoldFont=*-bold,ItalicFont=*-italic,BoldItalicFont=*-bolditalic]
\setmonofont{texgyrecursor}[Extension=.otf,UprightFont=*-regular,BoldFont=*-bold,ItalicFont=*-italic,BoldItalicFont=*-bolditalic,Scale=MatchLowercase]
\usepackage{amsmath,amssymb}
\usepackage{unicode-math}
\setmathfont{texgyrepagella-math.otf}
\usepackage{xcolor,microtype,enumitem,needspace,fancyhdr}
\usepackage{bookmark}
\definecolor{ink}{HTML}{17344B}
\definecolor{accent}{HTML}{197D80}
\definecolor{muted}{HTML}{596773}
\hypersetup{colorlinks=true,linkcolor=accent,urlcolor=accent,pdftitle={IE-24 - RILU
conditioning for the Neumann problem on smooth planar
domains},pdfauthor={Open Problems in Numerical Linear Algebra}}
\urlstyle{same}
\setlength{\parindent}{0pt}
\setlength{\parskip}{5pt plus 1pt minus 1pt}
\linespread{1.04}
\setlength{\emergencystretch}{3em}
\widowpenalty=10000
\clubpenalty=10000
\displaywidowpenalty=10000
\allowdisplaybreaks[1]
\setlist{nosep,leftmargin=1.5em}
\providecommand{\tightlist}{\setlength{\itemsep}{0pt}\setlength{\parskip}{0pt}}
\providecommand{\pandocbounded}[1]{#1}
\pagestyle{fancy}
\fancyhf{}
\fancyhead[L]{\sffamily\footnotesize\color{muted}OPEN PROBLEMS IN NUMERICAL LINEAR ALGEBRA}
\fancyhead[R]{\sffamily\bfseries\footnotesize\color{accent}IE-24}
\fancyfoot[L]{\parbox[b]{0.9\textwidth}{\sffamily\fontsize{8}{10}\selectfont\color{muted}Editorial ratings; a bounded search cannot certify that a problem remains unsolved.\\Literature check: 2026-09-10}}
\fancyfoot[R]{\sffamily\footnotesize\color{muted}\thepage}
\renewcommand{\headrulewidth}{0.3pt}
\renewcommand{\headrule}{\hbox to\headwidth{\color{accent}\leaders\hrule height \headrulewidth\hfill}}
\makeatletter
\renewcommand{\section}{\Needspace{5\baselineskip}\@startsection{section}{1}{0pt}{12pt plus 2pt minus 2pt}{4pt}{\sffamily\bfseries\large\color{ink}}}
\renewcommand{\subsection}{\Needspace{4\baselineskip}\@startsection{subsection}{2}{0pt}{9pt plus 1pt minus 1pt}{3pt}{\sffamily\bfseries\normalsize\color{ink}}}
\makeatother
\setcounter{secnumdepth}{0}
\begin{document}
{\sffamily\small\color{muted}Linear systems and elimination\par}
\vspace{3pt}
{\raggedright\hyphenpenalty=10000\exhyphenpenalty=10000\sffamily\bfseries\fontsize{21}{24}\selectfont\color{ink}RILU
conditioning for the Neumann problem on smooth planar domains\par}
\vspace{7pt}
{\sffamily\small\textbf{Difficulty:} challenging\quad\textbf{Importance:} interesting
to the community\par}
{\sffamily\small\bfseries\color{accent}Status: Open\par}
\vspace{4pt}
{\color{accent}\hrule height 0.8pt}
\vspace{6pt}
\textbf{Rating rationale:} Uniform cut-cell estimates on arbitrary
smooth domains are challenging beyond rectangles; effective Neumann
preconditioning has community importance.

\subsection{Statement}\label{statement}

Let \(\Omega\subset\mathbb R^2\) be a fixed bounded connected domain
with smooth boundary. For grid width \(h>0\), retain the points
\(p\in h\mathbb Z^2\) whose square control cells \(C_p=p+[-h/2,h/2]^2\)
intersect \(\Omega\). Order them lexicographically and identify them
with \(1,\ldots,N_h\). For cells sharing a face, put

\[
w_{pq}=\frac{\operatorname{length}(C_p\cap C_q\cap\Omega)}h,
\]

and set other weights to zero. Define the Purvis--Burkhalter Neumann
matrix by \(a_p=\sum_qw_{pq}\), \(A_{pp}=a_p\), and \(A_{pq}=-w_{pq}\)
for \(p\ne q\). Write \(L\) for its strictly lower triangular part.

For \(0<r<1\), the source's relaxed incomplete LU preconditioner uses
\(E=\operatorname{diag}(e_1,\ldots,e_{N_h})\), where \(e_1=a_1\) and, in
increasing order,

\[
e_p=a_p-\sum_{\substack{q<p\\w_{pq}>0}}
\frac{w_{pq}}{e_q}
\left(w_{pq}+(1-r)\sum_{\substack{s>q\\s\ne p}}w_{qs}\right),
\qquad M=(L+E)E^{-1}(E+L^T).
\]

This is Definition 4 written with numbered grid points; sums contain
only neighboring cells. Define \(\kappa_+(M^{-1}A)\) as the largest
divided by the smallest \textbf{positive} eigenvalue of
\(M^{-1/2}AM^{-1/2}\), excluding the Neumann zero eigenvalue.

\textbf{Conjecture.} There are constants
\(c_\Omega,K_\Omega,h_\Omega>0\), independent of \(h\), such that with
\(r=c_\Omega h^2\), for all \(0<h<h_\Omega\) the pivots are positive and

\[
\kappa_+(M^{-1}A)\le K_\Omega h^{-1}.
\]

Here \(h_\Omega\) is small enough that \(r<1\). This makes the source's
``some moderate constant'' explicit as an existence quantifier; no
numerical meaning is assigned to ``moderate.'' The authors recommend
\(c_\Omega=1\) in practice.

\subsection{Numerical significance}\label{numerical-significance}

This asks whether a specified sparse preconditioner gives the observed
improvement in conditioning for general curved Neumann boundaries. The
source proves the rectangular case, while its three-dimensional
experiments do not support the same RILU claim.

\newpage
\subsection{References and status
check}\label{references-and-status-check}

\begin{itemize}
\tightlist
\item
  B. Lee and C. Min, \emph{Optimal preconditioners on solving the
  Poisson equation with Neumann boundary conditions}, J. Comput. Phys.
  \textbf{433} (2021), 110189,
  \href{https://doi.org/10.1016/j.jcp.2021.110189}{DOI}.
  \href{https://math.ewha.ac.kr/~chohong/publications/article_49_copy.pdf}{Author
  text}: §2.1, equation (3); §3, Definition 4 and the conjecture,
  pp.~4--5; §5 gives the rectangular-domain result.
\item
  G. Hwang et al., \emph{Localized Estimation of Condition Numbers for
  MILU Preconditioners on a Graph}, SIAM J. Numer. Anal. \textbf{64}
  (2026), 1443--1477, \href{https://doi.org/10.1137/24M1722134}{DOI};
  \href{https://arxiv.org/abs/2501.00245}{available preprint}, §2.1 and
  §5, p.~21: assumes positive definite matrices and leaves Neumann
  conditions to future work.
\end{itemize}

On 2026-09-10, checked the full 2021 author text, the later available
MILU preprint, its 2026 publication record, and targeted
RILU/Neumann/title searches including 2025--2026. No general-domain
proof or counterexample was located. The 2026 journal full text was not
available; the later-source comparison uses its primary abstract and
author preprint.

\subsection{Independent audit ---
2026-09-10}\label{independent-audit-2026-09-10}

The recurrence and smooth-domain quantifiers were independently compared
with Lee--Min, Definition 4 and the §3 conjecture. The §5 proof treats
rectangles, whose corners exclude them from the displayed
smooth-boundary class; it therefore does not justify a
partial-resolution label for this exact target. The full author PDF and
the later MILU preprint were checked. The latter assumes positive
definiteness in §2.1 and leaves Neumann conditions to future work in §5.
Targeted subsequent searches found no smooth-domain resolution.
\end{document}
